\documentclass{discussion}

\usepackage{amsmath, amsthm, amssymb}
% \usepackage{xcolor, hyperref, cleveref}
\usepackage{mathtools}
\usepackage{derivative}
\usepackage{graphicx}
\usepackage{tabularx}
\usepackage{tasks}
\usepackage{mdframed}
\usepackage{commands}

\settasks{
    label-width = 1em,
    item-indent = 16pt,
}

\discussion{2}
\author{Sections B04 and B06}
\date{Fall 2024}

\begin{document}

\section{Basic convergence tests}

\begin{question}
    Match the series on the left with the appropriate convergence test and conclusion on the right, evaluating the series if asked.

    \begin{center}
        \setlength{\tabcolsep}{5em}
        \begin{tabularx}{\textwidth}{lX}
            \(\displaystyle \sum_{n=1}^\infty \frac{1}{n^{3/2}}\) & Diverges by the divergence test \\
            \(\displaystyle \sum_{n=1}^\infty \frac{1}{\sqrt{1 + n}}\) & \linespread{1.4}\selectfont  It's a convergent geometric series that converges to \rule{1.5cm}{0.15mm} \\
            \(\displaystyle \sum_{n=2}^\infty \frac{1}{n \ln n}\) & \linespread{1.4}\selectfont  It's a telescoping series that converges to \rule{1.5cm}{0.15mm} \\
            \(\displaystyle \sum_{n=2}^\infty \frac{1}{n (\ln n)^2}\) & Converges by the \(p\)-series Test \\
            \(\displaystyle \sum_{n=1}^\infty \frac{1}{\sqrt{n}}\) & Diverges by the \(p\)-series Test \\
            \(\displaystyle \sum_{n=1}^\infty \frac{3^n}{4^n}\) & Converges by the integral test \\
            \(\displaystyle \sum_{n=1}^\infty \frac{3}{n(n+2)}\) & Diverges by the integral test \\
        \end{tabularx}
    \end{center}
\end{question}

\begin{question}
    The following boxed claim and proof are incorrect.
    Identify the mistake(s).

    \begin{mdframed}
        \textbf{Claim:} \(\displaystyle \sum_{n=1}^\infty n^3 \eu^{-n} \leq 6\)

        \noindent
        \textbf{Proof:} The integral test tells us that
        \begin{equation*}
            \int_1^\infty x^3 \eu^{-x} \odif{x} \leq \sum_{n=1}^\infty n^3 \eu^{-n} \leq \int_0^\infty x^3 \eu^{-x} \odif{x}.
        \end{equation*}
        We evaluate the right integral to find that
        \begin{equation*}
            \int_0^\infty x^3 \eu^{-x} \odif{x} = \left[ (-x^3 - 3x^2 - 6x - 6) \eu^{-x} \right]_{0}^{\infty} = 6.
        \end{equation*}
        Hence, the inequality holds.
    \end{mdframed}
    Challenge: correctly prove that \(\sum_{n=1}^\infty n^3 \eu^{-n}\) converges and find an actual upper bound.
    (Check your answer with the fact that \(\sum_{n=1}^\infty n^3 \eu^{-n} \approx 6.0065\).)
\end{question}

\section{Estimating errors}

Let \(a_n \geq 0\) be a nonnegative sequence.
By estimating the infinite series \(\sum_{n=1}^\infty a_n\) with a partial sum \(\sum_{n=1}^{N} a_n\) to a precision of \(\epsilon\), we mean finding an integer \(N\) such that
\begin{equation*}
    \sum_{n=1}^\infty a_n - \sum_{n=1}^{N} a_n = \sum_{n=N+1}^\infty a_n < \epsilon.
\end{equation*}
This can be done by comparing the series to an integral.
If \(f\) is a decreasing function such that \(f(n) = a_n\), then integral test tells us that
\begin{equation*}
    \int_{N+1}^\infty f(x) \odif{x} \leq \sum_{n=N+1}^\infty a_n \leq \int_{N}^\infty f(x) \odif{x},
\end{equation*}
so the task becomes finding an integer \(N\) such that
\begin{equation*}
    \sum_{n=N+1}^\infty a_n \leq \int_{N}^\infty f(x) \odif{x} < \epsilon.
\end{equation*}

\begin{question}[Estimating errors]
    Estimate a value of \(N\) such that \(\sum_{n=1}^{N} a_n\) approximates the infinite series \(\sum_{n=1}^\infty a_n\) to within a precision of \(0.03\).
    \begin{tasks}(2)
        \task \(\displaystyle a_n = \frac{1}{n^2}\)
        \task \(\displaystyle a_n = \frac{1}{(n+1) (\ln (n+1))^2}\)
        \task* \(\displaystyle a_n = \frac{1}{n(n+1)}\) (Challenge: use another method besides comparison with an integral)
    \end{tasks}
\end{question}

\section{Comparison tests}

\begin{question}
    The following boxed claim and proof are incorrect.
    Identify the mistake(s).

    \begin{mdframed}
        \textbf{Claim:} \(\displaystyle \sum_{n=1}^{\infty} \frac{\sin(1/n)}{n}\) diverges.

        \noindent
        \textbf{Proof:}
        Since \(\sin(1/n) \leq 1\) for every positive integer \(n\), we can apply direct comparison test with the harmonic series \(\displaystyle \sum_{n=1}^{\infty} \frac{1}{n}\).
        Since the harmonic series diverges, the original series must also diverge.
    \end{mdframed}

    Challenge: show that \(\sum_{n=1}^{\infty} \frac{\sin(1/n)}{n}\) converges.
\end{question}

\begin{question}
    Match the series with the appropriate conclusion and state the series used to compare.

    \begin{center}
        \setlength{\tabcolsep}{5em}
        \begin{tabularx}{\textwidth}{lX}
            \(\displaystyle \sum_{n=1}^\infty \frac{n}{3^n}\) & \linespread{1.4}\selectfont Converges because it grows at the same rate as the convergent series: \\
            \(\displaystyle \sum_{n=1}^\infty \frac{n + 2}{n^3 - 2n + 2}\) & \linespread{1.4}\selectfont Diverges because it grows at the same rate as the divergent series: \\
            \(\displaystyle \sum_{n=1}^\infty \sin \left( \frac{1}{n} \right)\) & \linespread{1.4}\selectfont Converges because it grows slower than the convergent series: \\
            \(\displaystyle \sum_{n=1}^\infty \frac{1}{1 + \ln n}\) & \linespread{1.4}\selectfont Diverges because it grows faster than the divergent series:
        \end{tabularx}
    \end{center}
\end{question}

\end{document}